\section{Aplicación práctica: revisión de incidentes y recuperación de la confianza}

\subsection{Por qué la «revisión» es la segunda etapa de madurez del AI SRE}

Que el incidente se haya resuelto no significa que el trabajo haya terminado. Lo esencial está en la «revisión y la recuperación de la confianza»:

\begin{itemize}
    \item Registrar todas las interacciones del Agent, las llamadas a herramientas, las sugerencias automáticas y las decisiones humanas
    \item Analizar por qué la IA no logró dar la sugerencia correcta (falta de datos, inyección de prompt, hallucination del modelo)
    \item Reponer la documentación faltante, los hooks y las señales de monitoreo
    \item Comunicar al equipo las lecciones aprendidas y actualizar el runbook y las políticas de aprobación
\end{itemize}

\begin{warningbox}{Sin revisión no hay confianza}
La confianza que un equipo en operación deposita en la IA proviene de una sucesión de casos exitosos y de una ruta transparente para el manejo de errores. Si cada vez que la IA falla lo único que ocurre es «revertir y olvidar», el equipo no le permitirá asumir más permisos.
\end{warningbox}

\subsection{Flujo visual de la revisión de incidentes}

\begin{table}[H]
\centering
\begin{tabular}{p{0.2\textwidth} p{0.25\textwidth} p{0.28\textwidth}}
\toprule
Paso & Producto & Indicador clave \\
\midrule
Recopilar evidencia & Historial de comandos + prompt + tool outputs & Cobertura, integridad de los registros \\
Análisis de atribución & Cadena de causa raíz + matriz de evidencia & Margen de error, grado de correlación \\
Plan de remediación & Nuevo hook / nueva prueba / ajuste de permisos & Número de ejecuciones, número de Rollback \\
Consolidación del conocimiento & Actualización del runbook, etiquetado, reunión de difusión & Grado de completitud de la documentación, verificación de auditoría \\
\bottomrule
\end{tabular}
\caption{Flujo visual de la revisión de incidentes}
\end{table}

\subsection{Resumen de la sección}

El último tramo del manejo de incidentes es la revisión y la recuperación de la confianza. Registrar la cadena original de decisiones, analizar la causa raíz, reparar el monitoreo y convertir las lecciones en un runbook aplicable son la condición para que el AI SRE llegue realmente a producción.
